\setDocID{QA-V\&V-SPRINT1}
\setDocEstado{Release v1.0.0}

\chapter{Ciclo de Verificación y Validación para el Sprint 1}
\label{ch:sprint1}

El presente capítulo documenta de manera exhaustiva el primer ciclo de Verificación y Validación (V\&V) ejecutado para la plataforma EthosHub, correspondiente al Sprint 1. Elaborado bajo los lineamientos y rigor técnico del estándar IEEE 1012-2024, este informe consolida la estrategia de pruebas, la trazabilidad de los requisitos y los resultados de la ejecución técnica.\par
Durante esta iteración, el esfuerzo de Control de Calidad (QA) se concentró en auditar la estabilidad de la arquitectura fundacional del sistema. Esto abarca el motor de autenticación segura (incluyendo el registro tradicional y mediante proveedores externos OAuth2), la correcta asignación de roles de usuario, y el módulo de gestión de identidad, biografía y trayectoria profesional. A través de la ejecución sistemática de casos de prueba, el reporte documentado de defectos y su posterior remediación, este ciclo certifica la robustez de los módulos implementados y fundamenta el dictamen de calidad para su liberación a producción.

\newpage
\begin{NotaTecnica}
El objetivo de este sprint se centra en construir los cimientos funcionales de la plataforma EthosHub. Esto incluye el motor de acceso, la diferenciación de roles (profesionales y reclutadores), la seguridad de las cuentas y la gestión del perfil público (trayectoria laboral, certificaciones y redes sociales). Los resultados de este ciclo determinarán la viabilidad del pase a producción (ver, pag \pageref{sec:svvr}, sec \ref{sec:svvr}).
\end{NotaTecnica}

% ------------------------------------------
% 1. PLAN DE VERIFICACIÓN Y VALIDACIÓN
% ------------------------------------------
\section{Plan de Pruebas (Test Plan) - Sprint 1}
\label{sec:plan_pruebas}

\subsection{Alcance}
El presente plan de pruebas abarca la validación funcional de los módulos principales del sistema, con el objetivo de asegurar el correcto comportamiento de las funcionalidades más críticas. Dentro de este alcance se incluyen los procesos de registro de usuarios, tanto para profesionales como para reclutadores, así como el inicio y cierre de sesión. También se contempla la autenticación mediante proveedores externos, la recuperación de contraseñas y la gestión del perfil del usuario, incluyendo información como biografía, redes sociales y disponibilidad. La cobertura de estas funcionalidades se detalla en la Matriz de Trazabilidad (ver, pag \pageref{sec:rtm}, sec \ref{sec:rtm}).

\subsection{Objetivo}
El principal objetivo de este plan de pruebas es verificar que las funcionalidades implementadas cumplan con los criterios de aceptación definidos para cada historia de usuario. Asimismo, se busca validar que los flujos de autenticación, tanto tradicionales como externos, funcionen de manera correcta y sin interrupciones.

\subsection{Entorno y Herramientas}
\begin{itemize}[noitemsep]
    \item \textbf{Gestión de código:} Git / GitHub
    \item \textbf{Backend:} Java 17, Spring Boot 3.4.1
    \item \textbf{Frontend:} React + Vite + Tailwind + TypeScript
    \item \textbf{Herramientas de prueba:} DevTools (logs, network), Neon.tech (pruebas de BD)
    \item \textbf{Ambientes:} Entorno de desarrollo (local)
\end{itemize}

\subsection{Criterios de Entrada y Salida}
El proceso de pruebas se considera finalizado cuando todos los casos de prueba (ver, pag \pageref{sec:test_cases}, sec \ref{sec:test_cases}) han sido ejecutados. Asimismo, los bugs críticos y de alta prioridad (ver, pag \pageref{sec:bug_reports}, sec \ref{sec:bug_reports}) deben encontrarse resueltos o debidamente controlados. 

\textbf{Criterios de Ejecución Especiales:} Para el presente Sprint, se determina que \textbf{no se aplicó criterio de suspensión}, ya que los defectos encontrados no bloquearon el flujo principal de pruebas de manera catastrófica. Por consiguiente, \textbf{tampoco fue necesario aplicar un criterio de reanudación}.

% ------------------------------------------
% 2. MATRIZ DE TRAZABILIDAD (RTM)
% ------------------------------------------
\section{Matriz de Trazabilidad de Requisitos (RTM)}
\label{sec:rtm}

A continuación, se mapean las Historias de Usuario del Backlog del Sprint 1 con su estado de desarrollo y prioridad, base fundamental para la posterior ejecución de los Casos de Prueba (ver, pag \pageref{sec:test_cases}, sec \ref{sec:test_cases}).

\vspace{0.3cm}
\noindent
\fontsize{9}{10.8}\selectfont\sffamily
\begin{tabularx}{\textwidth}{|l|X|c|c|c|}
\hline
\rowcolor{headerTabla}
\textbf{ID} & \textbf{Historia de Usuario} & \textbf{Est. Esf.} & \textbf{Prioridad} & \textbf{Estado} \\
\hline
IDEN-001 & Configuración de Identidad Visual Crítica (Nombre, Titular y Foto). & 3 & Alta & \estadoPass \\
\hline
\rowcolor{grisTecnico}
LOG-001 & Registro de Usuario Profesional & 5 & Alta & \estadoPass \\
\hline
LOG-002 & Registro de Usuario Reclutador & 5 & Alta & \estadoPass \\
\hline
\rowcolor{grisTecnico}
LOG-003 & Inicio de Sesión de Usuarios (tradicional) & 3 & Alta & \estadoPass \\
\hline
LOG-005 & Recuperación de Contraseña & 8 & Alta & \estadoPass \\
\hline
\rowcolor{grisTecnico}
LOG-006 & Cierre de Sesión & 2 & Alta & \estadoPass \\
\hline
IDEN-003 & Narrativa Profesional: Editor de Biografía ("About Me") & 5 & Alta & \estadoPass \\
\hline
\rowcolor{grisTecnico}
IDEN-004 & Línea de Tiempo interactiva de Experiencia Laboral. & 5 & Alta & \estadoPass \\
\hline
IDEN-007 & Herramienta de Moderación y Auditoría (Admin Tool). & 8 & Alta & \estadoBlocked \\
\hline
\rowcolor{grisTecnico}
LOG-004 & Autenticación con Proveedores Externos & 5 & Media & \estadoPass \\
\hline
LOG-007 & Actualización de datos de perfil & 3 & Media & \estadoPass \\
\hline
\rowcolor{grisTecnico}
IDEN-008 & Optimización Técnica, SEO y Accesibilidad del Hero. & 3 & Media & \estadoBlocked \\
\hline
IDEN-009 & Gestión de Trayectoria Académica y Certificaciones. & 5 & Media & \estadoPass \\
\hline
\rowcolor{grisTecnico}
IDEN-002 & Gestión de Disponibilidad y Badge de Estado en tiempo real. & 3 & Media & \estadoPass \\
\hline
IDEN-005 & Agregador de Identidad Social (Social Links). & 3 & Media & \estadoPass \\
\hline
\rowcolor{grisTecnico}
IDEN-006 & Vista de Invitado y Barrera de Registro (Guest Teaser). & 5 & Media & \estadoPass \\
\hline
\end{tabularx}

% ------------------------------------------
% 3. CASOS DE PRUEBA
% ------------------------------------------
\clearpage
\section{Ejecución de Casos de Prueba}
\label{sec:test_cases}

El identificador de cada caso de prueba mapea directamente con la épica y la Historia de Usuario correspondiente definida en la Matriz de Trazabilidad (ver, pag \pageref{sec:rtm}, sec \ref{sec:rtm}).

\subsection{Módulo: Autenticación y Registro (LOG)}
\label{subsec:tc_log}

\subsubsection{Proceso de Registro de Usuarios (LOG-001 y LOG-002)}
\noindent
\fontsize{9}{10.8}\selectfont\sffamily
\begin{tabularx}{\textwidth}{|l|X|c|}
\hline
\rowcolor{headerTabla}
\textbf{ID Caso} & \textbf{Descripción de la Verificación} & \textbf{Estado} \\
\hline
TC-LOG-001-01 & Verificar que el sistema muestre un mensaje de confirmación al completar el registro. & \estadoPass \\ \hline
\rowcolor{grisTecnico}
TC-LOG-001-02 & Verificar que el sistema asigne correctamente el rol de profesional. & \estadoPass \\ \hline
TC-LOG-002-01 & Verificar que el sistema asigne correctamente el rol de reclutador al registrar empresa. & \estadoPass \\ \hline
\rowcolor{grisTecnico}
TC-LOG-001-03 & Verificar que el sistema muestre mensajes de error con caracteres no permitidos. & \estadoPass \\ \hline
TC-LOG-001-04 & Verificar que el sistema procese una sola vez la solicitud ante múltiples clics. & \estadoPass \\ \hline
\rowcolor{grisTecnico}
TC-LOG-001-05 & Verificar que el sistema permita el registro desde diferentes navegadores. & \estadoPass \\ \hline
TC-LOG-001-06 & Verificar almacenamiento correcto de datos de usuario en la BD tras el registro. & \estadoPass \\ \hline
\rowcolor{grisTecnico}
TC-LOG-001-07 & Verificar que se permita iniciar sesión posterior con credenciales registradas. & \estadoPass \\ \hline
TC-LOG-001-08 & Verificar visualización de error controlado ante falla interna en registro. & \estadoPass \\ \hline
\rowcolor{grisTecnico}
TC-LOG-001-09 & Verificar validación en tiempo real de campos obligatorios en formulario. & \estadoPass \\ \hline
\end{tabularx}

\subsubsection{Inicio y Cierre de Sesión (LOG-003 y LOG-006)}
\noindent
\fontsize{9}{10.8}\selectfont\sffamily
\begin{tabularx}{\textwidth}{|l|X|c|}
\hline
\rowcolor{headerTabla}
\textbf{ID Caso} & \textbf{Descripción de la Verificación} & \textbf{Estado} \\
\hline
TC-LOG-003-01 & Verificar inicio de sesión con credenciales válidas. & \estadoPass \\ \hline
\rowcolor{grisTecnico}
TC-LOG-003-02 & Verificar rechazo de acceso para correos no registrados. & \estadoPass \\ \hline
TC-LOG-003-03 & Verificar denegación de acceso por contraseña incorrecta. & \estadoPass \\ \hline
\rowcolor{grisTecnico}
TC-LOG-003-04 & Verificar mensajes de validación al iniciar sesión con campos vacíos. & \estadoPass \\ \hline
TC-LOG-003-05 & Verificar error por formato de correo inválido en login. & \estadoPass \\ \hline
\rowcolor{grisTecnico}
TC-LOG-003-06 & Verificar bloqueo de acceso para cuentas desactivadas. & \estadoPass \\ \hline
TC-LOG-003-07 & Verificar creación de sesión activa tras login exitoso. & \estadoPass \\ \hline
\rowcolor{grisTecnico}
TC-LOG-003-08 & Verificar redirección al panel correspondiente post-login. & \estadoPass \\ \hline
TC-LOG-003-09 & Verificar redirección al login ante intento de acceso a ruta protegida sin sesión. & \estadoPass \\ \hline
\rowcolor{grisTecnico}
TC-LOG-003-10 & Verificar mantenimiento de sesión durante la navegación. & \estadoPass \\ \hline
TC-LOG-006-01 & Verificar cierre correcto de sesión mediante la opción de logout. & \estadoPass \\ \hline
\rowcolor{grisTecnico}
TC-LOG-006-02 & Verificar bloqueo de acceso a rutas protegidas tras cerrar sesión. & \estadoPass \\ \hline
\end{tabularx}

\clearpage
\subsubsection{Autenticación Externa OAuth (LOG-004)}
\noindent
\fontsize{9}{10.8}\selectfont\sffamily
\begin{tabularx}{\textwidth}{|l|X|c|}
\hline
\rowcolor{headerTabla}
\textbf{ID Caso} & \textbf{Descripción de la Verificación} & \textbf{Estado} \\
\hline
TC-LOG-004-01 & Verificar visualización de botones de proveedores externos en login. & \estadoPass \\ \hline
\rowcolor{grisTecnico}
TC-LOG-004-02 & Verificar redirección al proveedor seleccionado para autenticación externa. & \estadoFail \\ \hline
TC-LOG-004-03 & Verificar creación automática de cuenta al autenticarse por primera vez. & \estadoFail \\ \hline
\rowcolor{grisTecnico}
TC-LOG-004-04 & Verificar prevención de cuentas duplicadas si ya existe asociación. & \estadoPass \\ \hline
TC-LOG-004-05 & Verificar retorno al login si el usuario cancela la autenticación externa. & \estadoPass \\ \hline
\rowcolor{grisTecnico}
TC-LOG-004-06 & Verificar rechazo de acceso si el proveedor no devuelve información válida. & \estadoPass \\ \hline
TC-LOG-004-07 & Verificar creación de sesión activa tras autenticación externa exitosa. & \estadoPass \\ \hline
\end{tabularx}

\subsubsection{Recuperación de Contraseña (LOG-005)}
\noindent
\fontsize{9}{10.8}\selectfont\sffamily
\begin{tabularx}{\textwidth}{|l|X|c|}
\hline
\rowcolor{headerTabla}
\textbf{ID Caso} & \textbf{Descripción de la Verificación} & \textbf{Estado} \\
\hline
TC-LOG-005-01 & Verificar envío de solicitud de recuperación con correo válido. & \estadoPass \\ \hline
\rowcolor{grisTecnico}
TC-LOG-005-02 & Verificar error al intentar recuperar con campo vacío. & \estadoPass \\ \hline
TC-LOG-005-03 & Verificar error de formato de correo en recuperación. & \estadoPass \\ \hline
\rowcolor{grisTecnico}
TC-LOG-005-04 & Verificar generación de token de recuperación válido. & \estadoPass \\ \hline
TC-LOG-005-05 & Verificar almacenamiento del token con tiempo de expiración. & \estadoPass \\ \hline
\rowcolor{grisTecnico}
TC-LOG-005-06 & Verificar validación de token no expirado. & \estadoPass \\ \hline
TC-LOG-005-07 & Verificar rechazo de restablecimiento con token inválido/inexistente. & \estadoPass \\ \hline
\rowcolor{grisTecnico}
TC-LOG-005-08 & Verificar rechazo de restablecimiento con token expirado. & \estadoPass \\ \hline
TC-LOG-005-09 & Verificar error si nueva contraseña y confirmación no coinciden. & \estadoPass \\ \hline
\rowcolor{grisTecnico}
TC-LOG-005-10 & Verificar actualización de contraseña e invalidación de token post-restablecimiento. & \estadoPass \\ \hline
\end{tabularx}

\subsubsection{Actualización de Perfil de Usuario (LOG-007)}
\noindent
\fontsize{9}{10.8}\selectfont\sffamily
\begin{tabularx}{\textwidth}{|l|X|c|}
\hline
\rowcolor{headerTabla}
\textbf{ID Caso} & \textbf{Descripción de la Verificación} & \textbf{Estado} \\
\hline
TC-LOG-007-01 & Verificar botón de guardado deshabilitado si hay errores en el formulario. & \estadoPass \\ \hline
\rowcolor{grisTecnico}
TC-LOG-007-02 & Verificar carga correcta de datos actuales al acceder a la edición de perfil. & \estadoPass \\ \hline
TC-LOG-007-03 & Verificar que el sistema permita modificar los campos del perfil. & \estadoPass \\ \hline
\rowcolor{grisTecnico}
TC-LOG-007-04 & Verificar bloqueo de guardado con errores de validación. & \estadoPass \\ \hline
TC-LOG-007-05 & Verificar actualización correcta de perfil con datos válidos. & \estadoPass \\ \hline
\rowcolor{grisTecnico}
TC-LOG-007-06 & Verificar almacenamiento en BD tras guardar el perfil. & \estadoPass \\ \hline
TC-LOG-007-07 & Verificar reflejo de cambios del perfil en otras vistas del sistema. & \estadoPass \\ \hline
\rowcolor{grisTecnico}
TC-LOG-007-08 & Verificar bloqueo de edición de perfil si no está autenticado. & \estadoPass \\ \hline
\end{tabularx}

\clearpage
\subsection{Módulo: Gestión de Identidad y Perfil (IDEN)}
\label{subsec:tc_iden}

\subsubsection{Identidad Visual, Disponibilidad y Biografía (IDEN-001, 002, 003)}
\noindent
\fontsize{9}{10.8}\selectfont\sffamily
\begin{tabularx}{\textwidth}{|l|X|c|}
\hline
\rowcolor{headerTabla}
\textbf{ID Caso} & \textbf{Descripción de la Verificación} & \textbf{Estado} \\
\hline
TC-IDEN-001-01 & Verificar visualización del nombre y titular en la cabecera del perfil. & \estadoPass \\ \hline
\rowcolor{grisTecnico}
TC-IDEN-001-02 & Verificar cambio de color de botón y copia al portapapeles en "Compartir". & \estadoPass \\ \hline
TC-IDEN-002-01 & Verificar renderizado de medalla de estado de disponibilidad. & \estadoPass \\ \hline
\rowcolor{grisTecnico}
TC-IDEN-003-01 & Verificar cambio a color rojo del contador con $\le$ 50 caracteres en biografía. & \estadoPass \\ \hline
TC-IDEN-003-02 & Verificar bloqueo de input al alcanzar límite de 500 caracteres. & \estadoPass \\ \hline
\rowcolor{grisTecnico}
TC-IDEN-001-03 & Verificar animación de carga y éxito al guardar el formulario sin errores. & \estadoPass \\ \hline
\end{tabularx}

\subsubsection{Redes Sociales y Enlaces (IDEN-005)}
\noindent
\fontsize{9}{10.8}\selectfont\sffamily
\begin{tabularx}{\textwidth}{|l|X|c|}
\hline
\rowcolor{headerTabla}
\textbf{ID Caso} & \textbf{Descripción de la Verificación} & \textbf{Estado} \\
\hline
TC-IDEN-005-01 & Verificar error por URL ajena a LinkedIn. & \estadoPass \\ \hline
\rowcolor{grisTecnico}
TC-IDEN-005-02 & Verificar aceptación de URL válida de GitHub. & \estadoPass \\ \hline
TC-IDEN-005-03 & Verificar bloqueo de guardado por errores en enlaces sociales. & \estadoPass \\ \hline
\rowcolor{grisTecnico}
TC-IDEN-005-04 & Verificar ausencia de errores si se dejan campos sociales vacíos. & \estadoPass \\ \hline
\end{tabularx}

\subsubsection{Gestión de Experiencia Laboral y Línea de Tiempo (IDEN-004, 006, 009)}
\noindent
\fontsize{9}{10.8}\selectfont\sffamily
\begin{tabularx}{\textwidth}{|l|X|c|}
\hline
\rowcolor{headerTabla}
\textbf{ID Caso} & \textbf{Descripción de la Verificación} & \textbf{Estado} \\
\hline
TC-IDEN-004-01 & Verificar error por fecha fin anterior a fecha inicio en experiencia. & \estadoPass \\ \hline
\rowcolor{grisTecnico}
TC-IDEN-004-02 & Verificar bloqueo del botón "Registrar Puesto" con error de fechas. & \estadoPass \\ \hline
TC-IDEN-004-03 & Verificar habilitación de botón al corregir rango de fechas. & \estadoPass \\ \hline
\rowcolor{grisTecnico}
TC-IDEN-004-04 & Verificar renderizado de línea vertical y puntos de cronología. & \estadoPass \\ \hline
TC-IDEN-004-05 & Verificar cálculo automático de duración en años y meses por puesto. & \estadoPass \\ \hline
\rowcolor{grisTecnico}
TC-IDEN-004-06 & Verificar deshabilitación de fecha fin al marcar "Estudiando aquí". & \estadoPass \\ \hline
TC-IDEN-004-07 & Verificar renderizado de la palabra "Presente" para estudios en curso. & \estadoPass \\ \hline
\rowcolor{grisTecnico}
TC-IDEN-004-08 & Verificar intercepción de mensaje de error desde API al fallar guardado. & \estadoPass \\ \hline
TC-IDEN-009-01 & Verificar estado "Lista vacía" en vista pública sin datos educativos. & \estadoPass \\ \hline
\rowcolor{grisTecnico}
TC-IDEN-004-09 & Verificar bloqueo de formulario si campos obligatorios nativos están vacíos. & \estadoPass \\ \hline
TC-IDEN-004-10 & Verificar deshabilitación del campo "Empresa" al marcar "Freelance". & \estadoPass \\ \hline
\rowcolor{grisTecnico}
TC-IDEN-004-11 & Verificar limpieza y bloqueo de fecha final al marcar trabajo actual. & \estadoPass \\ \hline
TC-IDEN-006-01 & Verificar requerimiento estricto de logo e imagen destacada. & \estadoPass \\ \hline
\rowcolor{grisTecnico}
TC-IDEN-006-02 & Verificar mensaje "Sin registros" al consultar base de datos vacía. & \estadoPass \\ \hline
TC-IDEN-004-12 & Verificar error por rango de fechas ilógico en Experiencia. & \estadoPass \\ \hline
\rowcolor{grisTecnico}
TC-IDEN-004-13 & Verificar renderizado del indicador verde animado para trabajo actual. & \estadoPass \\ \hline
TC-IDEN-004-14 & Verificar apertura de URL de empresa en nueva pestaña (seguridad). & \estadoPass \\ \hline
\rowcolor{grisTecnico}
TC-IDEN-004-15 & Verificar limpieza interna de fecha fin al cursar actualmente. & \estadoPass \\ \hline
TC-IDEN-006-03 & Verificar reemplazo de icono por "Freelance" en línea de tiempo. & \estadoPass \\ \hline
\end{tabularx}

\clearpage
\subsubsection{Herramientas Bloqueadas por Dependencia Externa (IDEN-007 y IDEN-008)}
\noindent
\fontsize{9}{10.8}\selectfont\sffamily
\begin{tabularx}{\textwidth}{|l|X|c|}
\hline
\rowcolor{headerTabla}
\textbf{ID Caso} & \textbf{Descripción de la Verificación} & \textbf{Estado} \\
\hline
TC-IDEN-007-01 & Verificar acceso al panel de administración para moderadores. & \estadoBlocked \\ \hline
\rowcolor{grisTecnico}
TC-IDEN-007-02 & Verificar carga de lista de usuarios reportados en la Admin Tool. & \estadoBlocked \\ \hline
TC-IDEN-007-03 & Verificar funcionalidad de baneo temporal de cuentas de usuario. & \estadoBlocked \\ \hline
\rowcolor{grisTecnico}
TC-IDEN-007-04 & Verificar registro de auditoría en la base de datos de moderación. & \estadoBlocked \\ \hline
TC-IDEN-007-05 & Verificar envío de correo automatizado tras sanción de cuenta. & \estadoBlocked \\ \hline
\rowcolor{grisTecnico}
TC-IDEN-007-06 & Verificar revocación manual de suspensión por parte de un SuperAdmin. & \estadoBlocked \\ \hline
TC-IDEN-007-07 & Verificar que cuentas suspendidas no puedan iniciar sesión en el portal. & \estadoBlocked \\ \hline
\rowcolor{grisTecnico}
TC-IDEN-008-01 & Verificar generación dinámica de meta-tags para perfiles públicos. & \estadoBlocked \\ \hline
TC-IDEN-008-02 & Verificar renderizado correcto de la vista previa en redes sociales (OpenGraph). & \estadoBlocked \\ \hline
\rowcolor{grisTecnico}
TC-IDEN-008-03 & Verificar navegación completa por teclado en el Hero del perfil. & \estadoBlocked \\ \hline
\end{tabularx}

% ------------------------------------------
% 4. REGISTRO DE DEFECTOS
% ------------------------------------------
\section{Registros de Defectos (Bug Reports)}
\label{sec:bug_reports}

\begin{AlertaDefecto}
Durante la ejecución de los Casos de Prueba (ver, pag \pageref{sec:test_cases}, sec \ref{sec:test_cases}), se detectaron los siguientes defectos. Cabe destacar que, para el cierre de este informe, todos los hallazgos aquí documentados fueron abordados y corregidos satisfactoriamente por el equipo de desarrollo.
\end{AlertaDefecto}

% BUG 001
\begin{BugReport}
ID de Defecto & \textbf{BUG-001}: El menú desplegable "Mi cuenta" no se despliega al hacer clic sobre el nombre de usuario. \\
\hline
\rowcolor{grisTecnico}
Severidad & Mayor \\
\hline
Prioridad & Media \\
\hline
\rowcolor{grisTecnico}
Pasos para Reproducir & 
\raggedright\arraybackslash
1. Iniciar sesión con una cuenta válida. \newline
2. Verificar sesión activa y nombre en navbar. \newline
3. Hacer clic sobre el nombre de usuario o avatar. \newline
4. Observar si aparece el menú de opciones. \\
\hline
Resultado Esperado & Se debe desplegar el menú emergente con: Mi perfil, Configuración y Cerrar sesión. \\
\hline
\rowcolor{grisTecnico}
Resultado Actual & \raggedright\arraybackslash No ocurre ninguna acción. El menú desplegable no aparece. \\
\hline
Ambiente & Chrome, Diseño Responsive. Asignado a: Mateo. \\
\hline
\end{BugReport}

% BUG 002
\clearpage
\begin{BugReport}
ID de Defecto & \textbf{BUG-002}: El mensaje legal de "Términos de Servicio" se mantiene fijo al hacer scroll en Registro. \\
\hline
\rowcolor{grisTecnico}
Severidad & Menor \\
\hline
Prioridad & Baja \\
\hline
\rowcolor{grisTecnico}
Pasos para Reproducir & 
\raggedright\arraybackslash
1. Navegar a la pantalla de "Crear Cuenta" (ver, pag \pageref{subsec:tc_log}, sec \ref{subsec:tc_log}). \newline
2. Ubicarse en el formulario. \newline
3. Hacer scroll hacia abajo. \newline
4. Observar el comportamiento del texto legal. \\
\hline
Resultado Esperado & El mensaje debe permanecer estático al final de la página (document flow). \\
\hline
\rowcolor{grisTecnico}
Resultado Actual & \raggedright\arraybackslash El mensaje actúa con CSS \comando{position: fixed}, superponiéndose al contenido al hacer scroll. \\
\hline
\end{BugReport}

% BUG 003
\begin{BugReport}
ID de Defecto & \textbf{BUG-003}: Inconsistencia en la validación del mínimo de caracteres en contraseña. \\
\hline
\rowcolor{grisTecnico}
Severidad & Moderada \\
\hline
Prioridad & Alta \\
\hline
\rowcolor{grisTecnico}
Pasos para Reproducir & 
\raggedright\arraybackslash
1. Acceder al formulario de registro. \newline
2. Ubicar el campo Contraseña. \newline
3. Observar el placeholder y el mensaje de error inferior. \\
\hline
Resultado Esperado & Información consistente (ej. "Mínimo 8 caracteres" en ambos lugares). \\
\hline
\rowcolor{grisTecnico}
Resultado Actual & \raggedright\arraybackslash Placeholder indica "Mínimo 6 caracteres", pero la validación exige "Mínimo 8 caracteres con mayúscula, minúscula, número y símbolo". \\
\hline
\end{BugReport}

% BUG 004
\begin{BugReport}
ID de Defecto & \textbf{BUG-004}: Inconsistencia en proveedores de autenticación externa. \\
\hline
\rowcolor{grisTecnico}
Severidad & Moderada \\
\hline
Prioridad & Media \\
\hline
\rowcolor{grisTecnico}
Pasos para Reproducir & 
\raggedright\arraybackslash
1. Acceder a inicio de sesión. \newline
2. Observar botones OAuth. \newline
3. Ir a Registro. \newline
4. Observar botones OAuth. \\
\hline
Resultado Esperado & Los mismos proveedores (ej. Google y GitHub) en ambas pantallas. \\
\hline
\rowcolor{grisTecnico}
Resultado Actual & \raggedright\arraybackslash Login muestra LinkedIn y GitHub. Registro muestra Google y GitHub. \\
\hline
\end{BugReport}

% BUG 005
\clearpage
\begin{BugReport}
ID de Defecto & \textbf{BUG-005}: Error 500 al registrarse con Google impide completar la autenticación. \\
\hline
\rowcolor{grisTecnico}
Severidad & \textcolor{red}{\textbf{Crítica}} \\
\hline
Prioridad & \textbf{Alta} \\
\hline
\rowcolor{grisTecnico}
Pasos para Reproducir & 
\raggedright\arraybackslash
1. Acceder a Registro. \newline
2. Clic en "Continuar con Google". \newline
3. Elegir cuenta válida. \newline
4. Continuar flujo. \\
\hline
Resultado Esperado & Autenticación exitosa y redirección al sistema. \\
\hline
\rowcolor{grisTecnico}
Resultado Actual & 
\raggedright\arraybackslash
\textit{Whitelabel Error Page} por excepción interna: \comando{No se encontro el atributo} \comando{OAuth2 requerido:} \comando{appAuthProvider} durante \comando{SuccessHandler}. \\
\hline
Evidencia & Ver traza de Spring Boot (Internal Server Error 500). \\
\hline
\end{BugReport}

% ------------------------------------------
% 5. INFORME DE RESUMEN
% ------------------------------------------
\section{Informe de Resumen de V\&V (SVVR)}
\label{sec:svvr}

\subsection{Evaluación Final de Calidad}
El ciclo de pruebas inicial del \textbf{Sprint 1} logró validar la infraestructura base de la plataforma. Si bien se detectaron incidencias críticas en la primera fase de ejecución, el equipo técnico abordó exitosamente las desviaciones levantadas. 

Por otro lado, existe un conjunto de Casos de Prueba marcados como Bloqueados (\estadoBlocked). Estos casos pertenecen a las Historias de Usuario \textbf{IDEN-007} y \textbf{IDEN-008}, las cuales fueron diferidas o puestas en pausa técnica temporal debido a dependencias en la arquitectura. Sin embargo, estas funcionalidades \textbf{no afectan el flujo base} de registro y gestión de perfiles del Sprint 1.

\subsection{Resumen de Defectos y Cobertura}
La cobertura de la ejecución refleja la totalidad de los Casos de Prueba documentados en la Matriz RTM:

\begin{itemize}[noitemsep]
    \item \textbf{Total Ejecutados (Exitosos):} 74 Casos de Prueba (\estadoPass).
    \item \textbf{Total Bloqueados (Postergados):} 10 Casos de Prueba (\estadoBlocked).
    \item \textbf{Defectos Abiertos (Fallos actuales):} 0 Casos de Prueba. Los 5 Bugs originales documentados (BUG-001 a BUG-005) se encuentran \textbf{Totalmente Resueltos}.
\end{itemize}

\subsection{Dictamen de Liberación}
\begin{tcolorbox}[colback=white, colframe=bgPass, boxrule=1.2pt, arc=2pt, title={\faCheckCircle\ \textbf{DICTAMEN: SPRINT EXITOSO Y APROBADO}}, coltitle=bgPass, fonttitle=\sffamily, attach title to upper=\par\vspace{0.2em}]
\sffamily
El Sprint 1 es declarado \textbf{Exitoso}. Se han validado las correcciones para todos los defectos levantados durante el ciclo. Dado que los Casos de Prueba bloqueados corresponden a herramientas accesorias (Admin Tool y Refinamiento SEO) que no comprometen la entrega de valor fundacional, el incremento del software resulta estable, robusto y funcional. Por lo tanto, \textbf{el producto es apto para su liberación a producción.}
\end{tcolorbox}

% ------------------------------------------
% 6. CONTROL DE VERSIONES
% ------------------------------------------
\clearpage
\section{Control de Versiones y Nomenclatura}
\label{sec:versioning}

Este documento y el incremento de software asociado han sido etiquetados con la versión \textbf{v1.0.0} siguiendo las directrices de \textit{Semantic Versioning (SemVer)}. 

\begin{itemize}
    \item \textbf{MAJOR (1):} Indica que se ha alcanzado el primer hito de producto estable y funcional. El sistema ya cuenta con el motor de acceso base y perfiles públicos, cumpliendo con la definición de "producto mínimo entregable" para este ciclo, lo que justifica el salto de la versión `0.x.x` (fase de desarrollo inicial) a la `1.x.x`.
    \item \textbf{MINOR (0):} Indica que, respecto a esta versión base de lanzamiento, aún no se han añadido nuevas características funcionales retrocompatibles (esto ocurrirá en el Sprint 2).
    \item \textbf{PATCH (0):} Representa que, habiendo subsanado todos los defectos levantados durante las pruebas antes del corte de la versión, no se han requerido parches de emergencia (\textit{hotfixes}) posteriores a la estabilización de este incremento.
\end{itemize}